\paragraph{擬似持ち上げのブロック周期}\label{par:false-lift-periods}
下降対の二つの空間を \(W,V\)、その階数を \(r_d,n\) とする。
\(Y\) を試験する形式の行列係数の複素共役とする。最高重み行列係数について、
調和多項式と有限ブロックデータでの下降が零である。
有限データを固定したまま無限部分を任意の Fock 多項式に取り替えると、
下降はこの零関数と \(U(W)\) の実点上の試験関数との畳み込みで表される。
従って多項式を取り替えた下降も零である。指標の変更が下降の変数の
行列式に掛ける因子であることから、これは本段で用いる規格化でも成り立つ。
\leanlinks{\leansrc{New/OneLift/FalseLifts/PeriodsOfFalseLifts}{調和データから任意の多項式への移行}}

有限部分の任意の試験関数は、レベルの素点 \(w\) での小さい剰余類の
指示関数に分けてよい。その中心を \(z_w=\binom{z'_w}{z''_w}\) とする。
全ての \(w\) で \(z'_w\) が可逆で \(z''_w(z'_w)^{-1}\) が整、かつ指定素点
ではこの比が剰余体で零なら、右から \((z'_w)^{-1}\) を掛けると
ブロックデータの値は \(1\) になる。十分小さい類では値は一定である。
指標付きでレベル群の平行移動を平均すると、その類の指示関数は
右移動したブロックデータの定数倍になる。これは、ブロックデータの
台上をレベル群が推移的に動かし、変換則が値を一点で決めるためである。
下降との対は平均で不変であり、既に零と分かったブロックデータの対に帰着する。

条件を満たさない素点があるときは、\(z'_w\) の核を使うか、
\(z''_w(z'_w)^{-1}\) の行に一様化元を掛けて原始な行を作る。
こうして \(lz_w=0\) となる整行 \(l\) で、指定行上の剰余
\(\lambda\) が非零のものを取れる。捻り素点では、指定行以外の
\(l\) の剰余は零にできる。有限体の可逆な \(s=1+v\lambda\) を
整行列 \(1+\widetilde v l\) に持ち上げる。これは \(z_w\) を固定するので
剰余類を保つ。他方、形式に対しては \(\chi_w(\det s)^{-1}\bar s\)
と作用する。指定素点では \(\bar s=s\) のブロック作用、捻り素点では
\(\bar s=1\) である。従って、これらの作用後の形式を各々同じ指示関数で
下降して足すと、元の下降値に群の位数を掛けた値になる。
指定素点では超平面和の仮定でこの和が零である。
捻り素点では \(\det(1+v\lambda)=1+\lambda v\) が剰余体の単元に
等しい重複度でわたり、\(\chi_w\ne1\) なので和が零になる。
こうしてレベル外で真空の全てのデータについて、\(W\) への下降が消える。
\leanlinks{\leansrc{New/OneLift/FalseLifts/BottomOfTower}{剰余類上の積分と超平面和}}

次に双曲平面を一枚ずつ加える。\(i\ge1\) 枚加えた階数を
\(m'=r_d+2i\)、\(k'=n-m'\) とし、下降を \(F\) とする。
混合模型で最後の等方直線の零項を取ると、その積分は \(i-1\) 枚での
下降なので零である。残りの核は一様に有界である。実際、標的が定値なので
その高さは \(1\) であり、始域の高さを \(H\) とすると、\(H>2\) では有理移動で
等方直線の長さを \(2/H\) 以下にできる。零項を引いた汎関数はその直線の
平均を消す。短い平行移動に対する任意次数の減衰を使うと、核の緩増大
\(O(H^A)\) を吸収できる。\(H\le2\) では緩増大の評価で足りる。
これをコンパクトな標的商で積分するので \(F\) は有界である。

レベル外の分解素点の剰余体位数を \(q\) とし、\(P_{d,w}(T)\) を
階数 \(d\) の Hecke 多項式とする。元の形式は
\(|X|=q^{-(n-1-r_d/2)}\) において \(P_{n,w}(X)\) で消える。
真空の局所 Hecke 作用を移すと \(F\) は球状で
\[
 \prod_{j=0}^{k'-1}(1-q^jX_0)\,P_{m',w}(X')F=0,\qquad
 X_0=\chi(\mathrm{Frob}_{\bar w})q^{m'/2}X,\qquad
 X'=\frac{q^{k'}X_0}{q^{n/2}} .
\]
ここで \(\chi\) は下降側の指標の基底変換である。
\(i\ge1\) と \(2r_d+2i+3<n\) より
\(q^{k'+m'/2}|X|\le1\) と \(q^{m'+1+k'/2}|X|\le1\) が成り立つ。
前者により積は非零、後者により \(q^{m'+1}|X'|\le1\) である。
\(P_{m',w}\) を剰余体上の部分空間の和に展開すると、零部分空間は恒等項で、
絶対係数の総和はガウス二項公式により
\[
 \prod_{j=0}^{m'-1}(1+q^j|X'|)
 \le\exp\!\left(\frac{1}{q(q-1)}\right)<2 .
\]
従って恒等項を残して移項すると \(\|F\|_\infty\) は
\(\|F\|_\infty\) の \(1\) 未満の倍数以下となり、\(F=0\) である。
最高重みベクトルの群軌道が係数空間を張るので、この消滅は全行列係数に及ぶ。
\leanlinks{\leansrc{New/OneLift/FalseLifts/HyperbolicPlaneStep}{零項を引いた核と小さい Hecke 根}}

\(V=V_1\oplus V_2\)、\(\dim V_1=N_d\)、\(2r_d+2N_d+3<n\) とする。
\(V_2\) 上の周期に現れる \(W\) の核を、\(V_1\) 上の核と掛けて \(U(V)\) 上の
核にする。\(V_1\) 側の試験データは、固定した周期の引数でその核が
非零となるように選べる。指標は二つのブロックで逆に取る。
この全体の核を \(Y\) に掛けた関数を \(Y_2\) とする。
\(Y_2\) の双曲空間への下降は、\(W\) と双曲空間の核の積の公式により、
いま零と分かった \(W\oplus\mathbb H^{N_d}\) への下降である。
有限部分の右移動もレベルの素点に支持される限りデータの条件を保つ。

\(I(g)\) を \(Y_2\) の \(U(V_2)\) 上のブロック周期を右移動したもの、
\(x(g)\) を対応する \(V_1\) の枠とする。混合模型で等方枠に沿う Fourier
係数を取ると、真空の Heisenberg 平行移動から
\(\psi_u(x(g))\,|e(\alpha_0 x(g))|^2\) が現れる。
\(\psi_u\) は成分と双対変数 \(u\) の加法指標であり、レベル外の真空に
固定した成分を除く。先に \(U(V_2)\) で積分すると、任意の \(u\) について
この重みと \(I(g)\) の積の積分が零となる。
真空の台では、枠の省いた成分の差をレベル外の整ユニタリ群で吸収できる。
従って周期の複素共役は加法指標の値を通じて連続関数として降りる。
指標の積と複素共役で生成する代数に Stone--Weierstrass の定理を適用すると
\[
 \int |e(\alpha_0x(g))|^2\,|I(g)|^2\,\mathrm dg=0 .
\]
被積分関数は連続かつ非負で、元の周期を与える枠で真空の重みは正であるから、
その周期は零である。\(V_1\) 側の非零の因子を除けば初めの \(V_2\) 周期の消滅を得る。
クラス関数に戻すとこの積分は安定化群で重み付けした類和である。
\leanlinks{\leansrc{Lib/Weil/Periods/DescentVanishing}{枠の Fourier 係数で周期を取り出す計算}
  \leanlinksep \leansrc{New/OneLift/FalseLifts/PeriodsOfFalseLifts}{部分群の類和での消滅}}
